\section{Reading long sources without losing their legal connections}
\label{sec:bounded-legal-context}

A sound interpretation may depend on passages that cannot all be presented
to a language model at once. A circular may state a duty in its main text,
qualify it in a footnote, and refer to a separate code for the relevant
standard. A judgment may define a relationship early in its reasons and
distinguish an apparently similar relationship much later. If a model
receives only the first passage, it can produce a fluent account of a rule
whose scope it has misunderstood. If every request receives the whole
collection, interpretation may never begin because the request exceeds
the application's capacity.

Long-document processing has established methods for this problem.
Overlapping chunks preserve nearby context; retrieval can restore a
passage's surroundings; hierarchical processing combines findings in stages;
and summary trees or graphs help locate distant connections. These methods
serve different purposes. Finding a few relevant passages is different from
reading every required source, and neither operation alone establishes that
all the legal relationships were understood. A useful design joins the
methods while retaining the original text against which conclusions can be
checked.

\subsection{Overlap preserves nearby context}

Consider a document represented by positions 0 to 100. Reading windows
0--50, 25--75 and 50--100 covers the whole document and repeats material
around the boundaries. A condition near position 49 can be read together
with a qualification near position 51. The overlap reduces the chance that
a local relation is severed merely because a fixed-size chunk ended there.

\begin{figure}[htbp]
\centering
\begin{tikzpicture}[x=0.105cm,y=0.6cm,font=\small]
  \draw[->] (0,0) -- (103,0);
  \foreach \x in {0,25,50,75,100} {
    \draw (\x,0.1)--(\x,-0.1);
    \node[below] at (\x,-0.1) {\x};
  }
  \draw[very thick,navy] (0,1)--(50,1);
  \draw[very thick,teal] (25,2)--(75,2);
  \draw[very thick,navy] (50,3)--(100,3);
  \node[right] at (101,1) {first};
  \node[right] at (101,2) {second};
  \node[right] at (101,3) {third};
  \draw[dashed] (10,0.2)--(10,1.2);
  \draw[dashed] (90,0.2)--(90,3.2);
  \node[above] at (10,1.2) {rule};
  \node[above] at (90,3.2) {exception};
\end{tikzpicture}
\caption{Every position is covered, but a rule at position 10 and an
exception at position 90 are never read together. A later comparison must
supply that connection. Positions are illustrative, not legal paragraph numbers.}
\label{fig:legal-overlap}
\end{figure}

The same example exposes overlap's limit. A rule at position 10 and an
exception at position 90 lie in different windows. There is no textual gap,
but an inference remains to be made. Increasing overlap helps local
continuity without guaranteeing that every consequential pair, or combination
of several passages, appears together. A later bounded review must bring
connected findings and their original supporting passages together.

OpenAI's public Python SDK documents a default maximum chunk size of 800
tokens and overlap of 400 tokens. Its static chunking declaration also
requires overlap not to exceed half the chunk size
\citep{openaichunking2026}. The ratio matches Figure~\ref{fig:legal-overlap}.
These are documented API parameters, not evidence that this width is optimal
for legal material. Paragraph and section boundaries should be preserved
where feasible, and quotations, definitions and exceptions near the chosen
boundaries need testing.

Tokens are units of model encoding, not bytes or characters. A request also
contains a question, instructions, source metadata and an output specification,
and needs room for the answer. A chunk that fills the nominal context window
leaves no space for these other requirements. Token accounting therefore needs
an explicit tokenizer and must remain separate from application byte limits.

\subsection{Retrieve a passage and restore its surroundings}

A short passage can be a good search result but a poor complete reading
context. Its distinguishing words may occur in one sentence, while the
subject of a pronoun or the scope of a condition appears nearby. Searching
and reading need not use identical units.

LlamaIndex's \texttt{SentenceWindowNodeParser} stores a sentence with
surrounding sentences as metadata. Its
\texttt{MetadataReplacementPostProcessor} can replace a retrieved sentence
with that window before the model reads it. Both operations matter: storing
a window without delivering it does not restore context. Parent-document
retrieval follows a similar principle. The inspected LangChain implementation
indexes small child chunks and returns larger parent sections or documents.
LlamaIndex's \texttt{AutoMergingRetriever} can merge retrieved children into
their parent and fill certain intervening gaps
\citep{llamaindexcontext2026,langchainparent2025}.

These operations restore actual text instead of asking a model to guess what
was omitted. They do not establish which distant exception applies. A parent
section can itself exceed the request budget, so expansion needs a measured
limit and an explicit unresolved result when necessary context cannot fit.
Identical words in two judgments also remain separate source occurrences;
text deduplication must not erase their different legal roles.

OpenAI's Responses File Search provides managed retrieval over uploaded
files. The inspected SDK exposes vector-store IDs, filters, result limits
and ranking controls, including weights for semantic and keyword matches
\citep{openaifilesearch2026}. A bounded set of relevant hits does not establish
that every required source was reviewed. These API details were verified in
public SDK code; the online guide endpoints were inaccessible during the
survey. No live service test is reported here.

\subsection{Read every piece and combine findings in stages}

Finding a relevant passage may be served by retrieval alone. Assessing every
duty in a circular requires a complete reading pass. Standard map/reduce
processing reads pieces separately and then combines their findings. If the
combined material still exceeds a request, reduction proceeds through
further bounded levels.

LlamaIndex's \texttt{TreeSummarize} implements this hierarchy. It repacks
text to the prompt budget, obtains responses for the pieces and recursively
combines those responses. It supports structured outputs as well as ordinary
text \citep{llamaindextree2026}. A legal application can reuse the packing
and combination mechanism while specifying the distinctions its output must
preserve.

A general summary can describe a document's main theme accurately while
omitting a rare exception that changes the present answer. A structured
reading should distinguish the proposed rule, conditions, exceptions,
actors, dates, supporting quotations, rival readings and unresolved
references. The original sources remain available separately. The record of
which sources and obligations were processed also stays outside the summaries;
summarizing that record would defeat its purpose.

Suppose the first window produces a duty and the third a possible exception.
Their findings can request a comparison containing both original passages,
the question and the necessary definitions. The comparison must explain
whether the exception qualifies this duty, concerns another actor or period,
or remains unresolved. Where a necessary dependency is unavailable or still
too large, the affected conclusion remains qualified while unrelated questions
continue.

This procedure moves the evidence needed for each inference into a bounded
request. It cannot make unlimited exact reasoning possible through arbitrary
compression. A reducer can omit findings, fail to shorten its input or even
expand it. Finite calls, explicit failure reporting and checks against original
sources are therefore necessary. A final request that reattaches the whole
raw dossier recreates the original capacity problem.

\subsection{RAPTOR connects passages through a hierarchy of topics}

Adjacent context is not always the most useful context.
\citet[sec.~3]{raptor2024} embed chunks, cluster related material, summarize
each cluster and repeat the process over the summaries. Their soft clustering
can associate a chunk with more than one topic. Related passages need not be
adjacent in the document.

RAPTOR retrieves from this hierarchy, either through its levels or from a
collapsed tree in which nodes at different levels compete for a bounded
context budget. The author implementation stores child identities alongside
summaries \citep{raptorcode2026}. A higher-level topic can consequently lead
back to underlying passages. That could help locate a definition or
qualification discussed in several places, while leaving the consequential
proposition to be checked against the original text.

The supporting evidence has a defined scope. Section~4 evaluates NarrativeQA,
QASPER and QuALITY. Appendices examine clustering, summarization, retrieval
and the contribution of different levels. Appendix~E reports hallucinations
in six of 150 inspected summaries and no discernible impact on the reported
question-answering results. That finding provides no assurance that a summary
preserves every legally consequential exception. The study motivates a
candidate retrieval method, not a legal acceptance verdict
\citep[sec.~4 and Apps.~B--F, I]{raptor2024}.

\subsection{GraphRAG connects entities and relationships}

Microsoft's GraphRAG extracts entities and relationships from chunks,
constructs hierarchical graph communities and produces community reports.
Its global search combines partial answers drawn from those reports. It
addresses questions whose answers are distributed across a collection
\citep[sec.~3]{graphrag2025}.

A legal investigation, for example, may need to distinguish provider,
intermediary, custodian and customer across documents. A graph can expose
proposed relationships for retrieval and comparison. But an extracted edge
such as “owns” already embodies an interpretation. Possession, legal title
and beneficial interest cannot be treated as interchangeable merely because
they connect the same names. The supporting text and uncertainty must survive
the extraction.

The retained paper evaluates global questions over podcast and news corpora
using model-judged measures including comprehensiveness and diversity. Its
statistical comparisons concern those tasks. Section~6.1 identifies domain
generalization and comparison of fabrication rates as further work
\citep[secs.~4--6 and App.~G]{graphrag2025}. These results do not establish
legal accuracy or exhaustive exception retention.

Implementation choices reinforce that distinction. The inspected global
reducer discards zero-scored points, ranks the remainder and takes what fits
its token budget. Malformed map output can become an empty result
\citep{graphragcode2026}. An exhaustive legal review must make such missing
coverage visible. The paper also evaluates source-text map/reduce, so a graph
is an alternative to test rather than a prerequisite for hierarchical reading.

\subsection{The design choice for LegalMath}

The justified starting point is an exact source store, overlapping pieces,
retrieval with restored context, complete structured reading and bounded
combination of connected findings. LlamaIndex supplies several components;
LangChain offers alternative parent retrieval. RAPTOR and GraphRAG can be
tested when simpler retrieval misses consequential distant connections.
OpenAI File Search is a managed retrieval alternative. Installing every
package would not establish which helps this product.

The capacity diagnosis gives the choice a concrete purpose. The partial
regulatory input exceeded an application ceiling of 200,000 bytes even after
exact sharing of repeated text. Later reviews repeatedly attached large
common evidence, so changing only the first chunk size could not solve the
overload. The proposed repair assigns sources and requirements to bounded
requests while retaining their identities across the complete reading.
The four original questions, 46 obligations and 35 retained readings remain
distinct targets; reducing their number is not a capacity repair.

A fair comparison uses the same sources and questions for ordinary overlapping
retrieval and the expanded-context method. Tests should include a distant
exception that changes the answer, a definition that changes the actor,
identical words in different sources, and a quotation or footnote attached to
the wrong proposition. Both delivered evidence and the resulting answer
matter. Returning the expected passages is an engineering result; drawing
and supporting the inference is a separate question.

The route remains optional until checks succeed on actual model responses.
Scripted local responses can establish source preservation, request bounds,
coverage accounting and failure handling, but cannot establish legal
interpretation quality. Retained papers, inspected code versions and source
hashes make the recommendation reproducible without implying completed
product acceptance.\footnote{The survey is retained in
\srcpath{docs/implementation/legal-mission-control/results/context-methods-2026-10-12.md};
source hashes are in
\srcpath{.localresources/legal-context-methods-2026-10-12/source-manifest.json}.
The revised plan is
\srcpath{docs/plans/legal-mission-capacity-resolution-2026-10-12.md}.}

